\chapter{Situating the Drivers on the Framework}
\label{ch:situating}

\chapterorient{We have met the six analyses (\cref{ch:targets}) and watched a
single run move from configuration to product through one lifecycle
(\cref{ch:lifecycle}). This chapter is where those threads tie together. We make
the shared \emph{skeleton} explicit, then ask the question the whole book turns
on: of the six analyses, how many genuinely different things happen in the
middle ``solve'' stage? The answer---\emph{four}---is the most important
taxonomy in the book. We name the four \emph{solve corners} and the three
\emph{reduction shapes}, and we lay out the grid that the rest of the book
fills in, one cell at a time. By the end you will be able to place any of the
six drivers on the framework and say exactly which stage it swaps.}

\section{The skeleton, made explicit}

Every analysis Palace runs, we claimed in \cref{ch:targets}, is the same
machine with a few parts swapped. \Cref{ch:lifecycle} showed the machine as an
actual program: parse the configuration, build a mesh, dispatch on the analysis
type, run a solve, and write out the product---with an optional adaptive
refinement loop wrapped around the solve. Strip that lifecycle down to its
load-bearing core and four words remain:
\[
  \texttt{config} \;\longrightarrow\; \texttt{assemble} \;\longrightarrow\;
  \texttt{solve} \;\longrightarrow\; \texttt{reduce} \;\longrightarrow\;
  \texttt{product}.
\]
This is \emph{the skeleton}. It is worth saying slowly, because everything in
the book hangs on it (\cref{fig:skeleton}).

\begin{itemize}[nosep]
  \item \textbf{config.} The user's description of the problem---geometry,
  materials, excitations, and the choice of analysis---parsed into a single
  read-only record. Nothing downstream ever mutates it.
  \item \textbf{assemble.} Turn the continuous physics on the mesh into one or
  more discrete \emph{operators}: large (conceptually) matrices like the
  stiffness $\Kop$, the mass $\Mop$, and the damping $\Cop$. Assembly is a
  \emph{fold} over the weak-form terms (\cref{ch:fem}); for now, treat it as the
  box that produces the operators.
  \item \textbf{solve.} Use those operators to compute the field, or fields, the
  analysis needs. This is the stage with all the variety, and the subject of
  the next section.
  \item \textbf{reduce.} Collapse the computed field(s) into the small physical
  quantity the user actually asked for---a capacitance matrix, a list of
  resonant frequencies, a table of S-parameters.
  \item \textbf{product.} The reduced quantity, written out. The reason the run
  existed.
\end{itemize}

\begin{figure}[t]
\centering
\begin{tikzpicture}[node distance=9mm and 11mm]
  \node[data] (cfg) {config};
  \node[stage, right=of cfg] (asm) {assemble};
  \node[stage, right=of asm] (solve) {\textbf{solve}};
  \node[stage, right=of solve] (red) {reduce};
  \node[product, right=of red] (prod) {product};
  \draw[flow] (cfg) -- (asm);
  \draw[flow] (asm) -- node[above,font=\scriptsize]{$\Kop,\Mop,\dots$} (solve);
  \draw[flow] (solve) -- node[above,font=\scriptsize]{field(s)} (red);
  \draw[flow] (red) -- (prod);
  % AMR back-edge
  \draw[backflow] (red.south) ++(0,-1mm)
    .. controls +(0,-7mm) and +(0,-7mm) .. (asm.south)
    node[midway, below, font=\scriptsize, text=simRust]
      {refine mesh (AMR outer fold)};
\end{tikzpicture}
\caption[The skeleton]{The skeleton every analysis shares:
\texttt{config} $\to$ \texttt{assemble} $\to$ \texttt{solve} $\to$
\texttt{reduce} $\to$ \texttt{product}.
The dashed back-edge is the adaptive-mesh-refinement outer loop
(\cref{ch:amr}): when refinement is on, the reduced error estimate feeds a mesh
update and the whole assemble--solve--reduce body runs again; when it is off,
the body runs exactly once. The bold \texttt{solve} box is where the six
analyses diverge---and they diverge into only four shapes.}
\label{fig:skeleton}
\end{figure}

One feature of the lifecycle does not fit inside a single left-to-right pass:
adaptive mesh refinement. When refinement is enabled, the run does not stop
after one reduce. Instead the error it estimates feeds a mesh update, and the
entire assemble--solve--reduce body runs again on the finer mesh---and again,
until the error is small enough or a step budget is spent. That is the dashed
back-edge in \cref{fig:skeleton}. It is an \emph{outer fold} wrapped around the
skeleton: each pass begins from the mesh the previous pass refined. When
refinement is off, the fold degenerates to a single pass, and the skeleton is
exactly the straight line. We will meet this outer fold again as a genuine
computation in \cref{ch:amr}; for now, note only that it leaves the
five-stage skeleton intact and merely runs it more than once.

\begin{keyidea}
The skeleton is \texttt{config} $\to$ \texttt{assemble} $\to$ \texttt{solve}
$\to$ \texttt{reduce} $\to$ \texttt{product},
optionally wrapped in an adaptive-refinement outer fold. The configuration is
read-only; assembly produces operators; the solve produces fields; the
reduction produces the physical product. Of the five stages, four are nearly
identical across all six analyses. \emph{Almost the entire difference between
the six analyses lives in the \texttt{solve} stage}---and, to a lesser degree,
in the \texttt{reduce} stage. This chapter classifies both.
\end{keyidea}

\section{The four solve corners}

Here is the central observation of the book. Look at what the ``solve'' stage
actually does in each of the six analyses, and ignore the physics: ignore which
equation, which function space, which materials. Ask only about the
\emph{shape} of the computation---how many solves happen, whether they depend on
one another, whether the operator stays fixed or changes, and whether we write
the loop or hand the whole problem to a library. When you strip the physics
away like this, the six analyses collapse onto just \emph{four} distinct shapes.
We call them the four \emph{solve corners}.

Two questions separate them (\cref{fig:corners}). First: \emph{does the operator
stay fixed, or does it change as we go?} An analysis might assemble its operator
once and reuse it for every solve, or it might have to rebuild the operator
before each solve. Second: \emph{is the computation a map or a fold?} A
\emph{map} runs a collection of \emph{independent} solves---each one could be
done on a separate computer, in any order, because none depends on another's
result. A \emph{fold} runs a \emph{chain} of solves---each one starts from the
previous one's output, so they are intrinsically sequential. Crossing these two
questions gives a $2\times 2$ grid; three of its cells are occupied by real
corners, and a fourth corner sits outside the grid entirely, because in it we do
not write a loop at all. We take the four in turn.

\begin{figure}[p]
\centering
\begin{tikzpicture}[font=\footnotesize]
  % axis labels
  \node[cornertag, anchor=south] at (0,4.55) {operator \emph{fixed}};
  \node[cornertag, anchor=south] at (4.6,4.55) {operator \emph{varying}};
  \node[cornertag, rotate=90, anchor=south] at (-2.55,3.3) {a \textbf{map} (independent)};
  \node[cornertag, rotate=90, anchor=south] at (-2.55,0.9) {a \textbf{fold} (chained)};

  % --- top-left: fixed-operator map ---
  \begin{scope}[shift={(-1.3,2.5)}]
    \node[stage, minimum width=24mm, minimum height=22mm, anchor=south west,
      fill=simBgBlue, align=center] (q1) at (0,0) {};
    \node[anchor=north, font=\scriptsize\bfseries] at (1.2,2.0)
      {Fixed-operator map};
    % one operator box, three independent solves off it
    \node[opbox, minimum width=7mm] (op1) at (0.35,1.0) {$\Kop$};
    \foreach \i/\y in {1/1.45,2/1.0,3/0.55}{
      \node[product, minimum width=4mm, minimum height=3.5mm, scale=0.6]
        (s1\i) at (1.95,\y) {};
      \draw[flow, shorten >=1pt] (op1.east) -- (s1\i.west);}
    \node[anchor=north, font=\scriptsize\ttfamily] at (1.2,0.42)
      {solve\_family};
    \node[anchor=north, font=\scriptsize\itshape] at (1.2,0.08)
      {electrostatic, magnetostatic};
  \end{scope}

  % --- top-right: operator-varying map ---
  \begin{scope}[shift={(3.3,2.5)}]
    \node[stage, minimum width=24mm, minimum height=22mm, anchor=south west,
      fill=simBgBlue, align=center] (q2) at (0,0) {};
    \node[anchor=north, font=\scriptsize\bfseries] at (1.2,2.0)
      {Operator-varying map};
    \foreach \i/\y in {1/1.5,2/1.0,3/0.5}{
      \node[opbox, minimum width=6mm, scale=0.75] (op2\i) at (0.5,\y)
        {$\Kop(\omega)$};
      \node[product, minimum width=4mm, minimum height=3.5mm, scale=0.6]
        (t2\i) at (1.95,\y) {};
      \draw[flow, shorten >=1pt] (op2\i.east) -- (t2\i.west);}
    \node[anchor=north, font=\scriptsize\ttfamily] at (1.2,0.42)
      {frequency\_sweep};
    \node[anchor=north, font=\scriptsize\itshape] at (1.2,0.08)
      {driven};
  \end{scope}

  % --- bottom-left: state-threaded fold ---
  \begin{scope}[shift={(-1.3,0.0)}]
    \node[stage, minimum width=24mm, minimum height=22mm, anchor=south west,
      fill=simBgGreen, align=center] (q3) at (0,0) {};
    \node[anchor=north, font=\scriptsize\bfseries] at (1.2,2.0)
      {State-threaded fold};
    \foreach \i/\x in {1/0.3,2/0.95,3/1.6}{
      \node[opbox, minimum width=4mm, minimum height=4mm, scale=0.7]
        (f\i) at (\x,1.0) {};}
    \draw[flow, shorten >=1pt] (f1.east) -- (f2.west);
    \draw[flow, shorten >=1pt] (f2.east) -- (f3.west);
    \node[anchor=north, font=\scriptsize\ttfamily] at (1.2,0.42)
      {fold\_solve};
    \node[anchor=north, font=\scriptsize\itshape] at (1.2,0.08)
      {transient, AMR loop};
  \end{scope}

  % --- bottom-right: black-box (outside the grid) ---
  \begin{scope}[shift={(3.3,0.0)}]
    \node[extern, minimum width=24mm, minimum height=22mm, anchor=south west,
      align=center] (q4) at (0,0) {};
    \node[anchor=north, font=\scriptsize\bfseries] at (1.2,2.0)
      {Opaque black-box};
    \node[extern, minimum width=12mm, minimum height=8mm] (bb) at (1.2,1.0)
      {\scriptsize library};
    \node[anchor=north, font=\scriptsize\ttfamily] at (1.2,0.42)
      {eigsolve};
    \node[anchor=north, font=\scriptsize\itshape] at (1.2,0.08)
      {eigenmode, boundary-mode};
  \end{scope}

  \node[font=\scriptsize\itshape, text=simViolet, align=center]
    at (4.5,-0.65) {(no loop we write)};
\end{tikzpicture}
\caption[The four solve corners]{The four solve corners. Two questions
classify a solve: is the operator \emph{fixed} or \emph{varying} as we proceed
(columns), and is the computation a \emph{map} of independent solves or a
\emph{fold} of chained solves (rows)? Three corners occupy the resulting grid.
A fixed-operator map (top-left) reuses one operator across independent solves; an
operator-varying map (top-right) rebuilds the operator before each independent
solve; a state-threaded fold (bottom-left) chains solves so each starts from the
last one's output. The fourth corner (bottom-right, dashed) sits \emph{outside}
the grid: we write no loop at all---we hand the whole problem to a library
routine. Each corner names its combinator and the drivers that use it.}
\label{fig:corners}
\end{figure}

\subsection{Corner 1: the fixed-operator map}

The simplest corner. Assemble the operator \emph{once}, then run a family of
independent solves that all reuse it. Each solve has the same operator on the
left-hand side and a different right-hand side; the solves do not depend on one
another, so they form a \emph{map} over the family of right-hand sides.

The exemplar is electrostatics (\cref{ch:electrostatics}). To compute the
capacitance of $n$ conductors, we assemble the stiffness operator $\Kop$ once,
then solve $\Kop\,\vx_i = \vb_i$ for $n$ different right-hand sides---one per
conductor held at unit voltage. The operator $\Kop$ never changes; only the
right-hand side does. We capture $\Kop$ once, outside the loop, and map the
solver over the family. Magnetostatics works identically, with a different
operator and a different family. In the framework's vocabulary this corner is
the combinator \code{solve\_family}, written
\[
  \texttt{solve\_family}\;\Kop\;[\vb_1,\dots,\vb_n]
    \;=\; [\,\texttt{solve}\;\Kop\;\vb_1,\;\dots,\;\texttt{solve}\;\Kop\;\vb_n\,].
\]
The single load-bearing fact is the \emph{hoist}: the operator is set up once,
before the map, not once per element. Everything else about this corner follows
from that.

\subsection{Corner 2: the operator-varying map}

Now keep the map---a family of independent solves---but change one thing: the
operator is no longer fixed. It must be rebuilt before each solve.

The exemplar is the driven analysis (\cref{ch:driven}). To compute S-parameters
across a frequency band, we solve the time-harmonic system at each frequency
$\omega$ in the sweep. But the system operator depends on frequency:
$\mat{A}(\omega) = \Kop + i\omega\,\Cop - \omega^2\Mop$. So inside the sweep,
\emph{before each solve}, we must reassemble $\mat{A}(\omega)$ at the current
frequency. The solves are still independent---frequency $\omega_3$ does not need
the answer at $\omega_2$---so this is still a map, not a fold. What distinguishes
it from corner~1 is precisely that the operator is rebuilt \emph{inside} the
loop rather than hoisted outside it. The combinator is \code{frequency\_sweep}.

This contrast---hoist the operator versus rebuild it---is small to state and
large in consequence, and it is the entire structural difference between
electrostatics and the driven analysis. We will return to it in detail in
\cref{ch:driven}; for now, hold onto the picture in \cref{fig:corners}: one
operator feeding many solves (corner~1) versus a fresh operator for each solve
(corner~2).

\begin{pitfall}
It is tempting to call corner~2 a ``harder version of corner~1'' and leave it
at that. It is harder, but the difference is not a matter of degree---it is a
matter of \emph{kind}. In corner~1 the expensive operator setup happens once and
is amortized over the whole family; in corner~2 it happens every single
iteration. An implementation that ``optimizes'' corner~2 by accidentally
hoisting the operator out of the loop does not run faster---it computes the
wrong answer, because it solves every frequency against the operator of the
first one. The position of one setup call relative to one loop is load-bearing.
\end{pitfall}

\subsection{Corner 3: the state-threaded fold}

Drop the independence. In corners~1 and~2 the solves were a map: independent,
reorderable, parallelizable. In corner~3 they are a \emph{fold}: each solve
begins from the state the previous solve produced. The computation is a chain,
and the chain cannot be reordered or parallelized, because step $k$ literally
needs the output of step $k-1$ as its input.

The exemplar is the transient analysis (\cref{ch:transient}). To compute the
time-domain response, we march forward in time: the field state at time
$t_{k+1}$ is computed from the field state at time $t_k$ by taking one time
step. The state is \emph{threaded} through the computation, each step
transforming it and handing it to the next. This is a \emph{fold}, the
fundamental sequential combinator we will study at length in \cref{ch:fold}; in
the framework it is \code{fold\_solve}. The adaptive-refinement outer loop from
\cref{fig:skeleton} is also a fold of this shape---each refinement pass begins
from the mesh the last pass produced---which is why the same combinator appears
for both.

The distinction between corner~3 and the two map corners is the distinction
between a \emph{movie} and a \emph{contact sheet}. A map produces a set of
independent snapshots that happen to be collected together; a fold produces a
single evolving history in which each frame depends on the last. We will see in
\cref{ch:fold} that the map is in fact the degenerate special case of the fold
in which each step ignores the running state---but the operational difference,
parallelizable versus strictly sequential, is real and central.

\subsection{Corner 4: the opaque black-box}

The fourth corner is different in nature from the other three. In corners~1--3
\emph{we} write the loop: we control the map or the fold, and we can see every
solve it performs. In corner~4 we write no loop at all. We assemble the
operators, hand the entire problem to a specialized library routine, and receive
the answer. The iteration that produces it happens \emph{inside} the library,
where our framework cannot see it.

The exemplar is the eigenmode analysis (\cref{ch:eigenvalues}). To find the
resonant frequencies of a cavity, we assemble the operator pencil
$(\Kop,\Cop,\Mop)$ once and pass it to an eigenvalue solver---SLEPc or ARPACK,
in Palace's case---with a single call. The solver runs its own sophisticated
iteration (a Krylov--Schur or Arnoldi process; \cref{ch:eigenvalues}) and
returns the converged eigenpairs. We do not see the iteration; we see one call
and its result. The boundary-mode analysis uses the same corner, on a
two-dimensional cross-section rather than a volume. In the framework this corner
is the combinator \code{eigsolve}, and it is deliberately \emph{opaque}: it
names the library boundary and the contract across it, and stops there.

This opacity is not a gap in our understanding; it is an honest record of where
Palace's authors chose to call out to someone else's code. \Cref{ch:eigenvalues}
both documents that boundary and---separately---reconstructs what the library
does internally, in our own vocabulary, so that the black box becomes a glass
box for the reader who wants to look inside.

\begin{keyidea}
The four solve corners are the most important taxonomy in this book. They are
distinguished by two questions plus one exception:
\begin{itemize}[nosep]
  \item \textbf{Fixed-operator map} (\code{solve\_family}): assemble once, then
  map an independent family of solves over the fixed operator. \emph{Hoist the
  operator.} Exemplars: electrostatic, magnetostatic.
  \item \textbf{Operator-varying map} (\code{frequency\_sweep}): a map of
  independent solves, but rebuild the operator \emph{inside} the loop because it
  changes each step. \emph{Do not hoist.} Exemplar: driven.
  \item \textbf{State-threaded fold} (\code{fold\_solve}): a chain, not a map---
  each step's input is the previous step's output. \emph{Sequential.}
  Exemplars: transient, and the AMR outer loop.
  \item \textbf{Opaque black-box} (\code{eigsolve}): no loop we write at all;
  one call to a library solver returns the answer. Exemplars: eigenmode,
  boundary-mode.
\end{itemize}
Every later chapter that studies a driver is, in part, the study of which corner
that driver sits in and why.
\end{keyidea}

\section{The three reduction shapes}

The solve stage has four shapes; the reduce stage has three. After a driver has
produced its field or family of fields, the \texttt{reduce} stage collapses
them into the small physical quantity the user wants. As with the solve stage,
the physics differs from analysis to analysis but the \emph{shape} of the
reduction does not---and there are only three shapes (\cref{fig:reductions}).

\begin{figure}[t]
\centering
\begin{tikzpicture}[font=\footnotesize]
  % --- (a) symmetric Gram matrix ---
  \begin{scope}[shift={(0,0)}]
    \foreach \i in {0,1,2}{
      \foreach \j in {0,1,2}{
        \pgfmathtruncatemacro{\sym}{(\i==\j)?1:0}
        \node[draw=simGray, semithick,
          fill={\ifnum\i=\j simBgRust\else simBgBlue\fi},
          minimum size=6.5mm, inner sep=0pt]
          (g\i\j) at (\j*0.66, -\i*0.66) {};}}
    % symmetry mirror line
    \draw[densely dashed, simRust, semithick]
      (-0.45,0.45) -- (1.77,-1.77);
    \node[font=\scriptsize\bfseries, align=center] at (0.66,-2.45)
      {Rank-2 Gram};
    \node[font=\scriptsize, align=center] at (0.66,-2.9)
      {symmetric: $G_{ij}=G_{ji}$};
    \node[font=\scriptsize\itshape, align=center] at (0.66,-3.4)
      {capacitance,\\inductance};
  \end{scope}

  % --- (b) non-symmetric port-projection ---
  \begin{scope}[shift={(4.2,0)}]
    \foreach \i in {0,1,2}{
      \foreach \j in {0,1,2}{
        \node[draw=simGray, semithick, fill=simBgGreen,
          minimum size=6.5mm, inner sep=0pt]
          (p\i\j) at (\j*0.66, -\i*0.66) {};}}
    \node[font=\scriptsize, text=simGreen] at (0.66,0.55)
      {ports out};
    \node[font=\scriptsize, rotate=90, text=simGreen] at (-0.62,-0.66)
      {ports in};
    \node[font=\scriptsize\bfseries, align=center] at (0.66,-2.45)
      {Rank-2 projection};
    \node[font=\scriptsize, align=center] at (0.66,-2.9)
      {\emph{not} symmetric};
    \node[font=\scriptsize\itshape, align=center] at (0.66,-3.4)
      {S-parameters};
  \end{scope}

  % --- (c) rank-1 scalar / mode table ---
  \begin{scope}[shift={(8.4,0)}]
    \foreach \i in {0,1,2,3}{
      \node[draw=simGray, semithick, fill=simBgGold,
        minimum width=18mm, minimum height=5mm, inner sep=0pt]
        (r\i) at (0, -\i*0.6) {};}
    \node[font=\scriptsize] at (0,0) {mode 1: $f_1, Q_1$};
    \node[font=\scriptsize] at (0,-0.6) {mode 2: $f_2, Q_2$};
    \node[font=\scriptsize] at (0,-1.2) {mode 3: $f_3, Q_3$};
    \node[font=\scriptsize] at (0,-1.8) {$\vdots$};
    \node[font=\scriptsize\bfseries, align=center] at (0,-2.45)
      {Rank-1 table};
    \node[font=\scriptsize, align=center] at (0,-2.9)
      {one row per element};
    \node[font=\scriptsize\itshape, align=center] at (0,-3.4)
      {eigenfreq.\,$+\,Q$,\\energy, modes};
  \end{scope}
\end{tikzpicture}
\caption[The three reduction shapes]{The three reduction shapes. A
\emph{rank-2 Gram} (left) is a symmetric matrix of energy-like inner products
over a family of solutions, $G_{ij}=w(i,j)\,\vx_j^{\!\top}\Kop\,\vx_i$; its
symmetry (dashed mirror) is structural. A \emph{rank-2 projection} (centre) is
also a matrix, but it projects an ``in'' family onto an ``out'' family and is
\emph{deliberately not} a Gram---it is generally non-symmetric. A \emph{rank-1
table} (right) is a one-dimensional list: one row of scalars per mode, per
domain, or per propagation mode. The distinction between the left and centre
shapes is the single most over-unified point in the design---see the pitfall.}
\label{fig:reductions}
\end{figure}

\subsection{The rank-2 symmetric Gram}

The first shape is a \emph{Gram matrix}: a symmetric matrix of energy-like inner
products formed over a family of solutions. Given a solution family
$\vx_1,\dots,\vx_n$ and the assembled operator $\Kop$, the Gram reduction
computes
\[
  G_{ij} \;=\; w(i,j)\;\vx_j^{\!\top}\,\Kop\,\vx_i,
\]
a matrix whose $(i,j)$ entry is an operator-weighted inner product of solutions
$i$ and $j$. Because $\Kop$ is symmetric, $G_{ij}=G_{ji}$: the result is a
symmetric matrix, structurally, not by coincidence.

Both capacitance and inductance are Gram reductions. They use the \emph{same
verb}---\code{gram\_reduce} in the framework---and differ only in the weight
$w(i,j)$. Capacitance uses the unit weight $w(i,j)=1$ (the energy form
$C_{ii}=2U/V^2$ with unit applied voltage); inductance uses the
current-normalized weight $w(i,j)=1/(I_i I_j)$. One verb, two weights, two
physical products. This is the cleanest instance of the book's recurring
refrain: the same machine with one part---here, a scalar weight function---
swapped. We study it fully in \cref{ch:reductions}.

\subsection{The rank-2 port-projection}

The second shape is also a matrix, and it is tempting to file it under the
first. Resist. The S-parameter matrix is a \emph{projection}, not a Gram. Its
entry $S_{ij}(\omega)$ records how much of a wave entering port $j$ leaves port
$i$: it projects the solved field for each driven port onto each output port's
mode. It is a matrix indexed by (out-port, in-port), but the two index families
play different roles, the operation is a linear projection rather than an
energy-weighted self-inner-product, and the result is in general
\emph{not symmetric}. The framework gives it its own verb,
\code{sparameter\_reduce}, deliberately distinct from \code{gram\_reduce}.

\begin{pitfall}
Two reductions producing matrices is not enough to make them the same reduction.
The Gram products (capacitance, inductance) and the port-projection
(S-parameters) both yield rank-2 results, and a tidy-minded designer is sorely
tempted to unify them under one ``matrix reduction'' verb. This is a mistake.
The Gram is a \emph{symmetric self-pairing} of a family with itself through an
operator; the port-projection is an \emph{asymmetric projection} of one family
onto a different basis. Forcing them into one verb either smuggles a false
symmetry into the S-matrix or burdens the Gram with machinery it does not need.
The right move---which the framework makes---is two verbs that happen to share
an output rank. \emph{Sharing a shape is not sharing a meaning.}
\end{pitfall}

\subsection{The rank-1 scalar or mode table}

The third shape is a \emph{table}: a one-dimensional list with one row per
element, each row holding a few scalars (or, in one case, a few small fields).
There is no matrix structure, no pairing of a family with itself---just a map
over a collection, producing one tuple per item.

Three products live here. Eigenfrequencies and quality factors form a table with
one row per mode: $(f_m, Q_m)$ for each resonant mode. Energy and participation
factors form a table with one row per material domain: how much of the field
energy sits in each region. And waveguide propagation modes form a table with
one row per mode, each carrying a propagation constant $k_n$ and a small
transverse field pattern (the one case where the row holds fields, not just
scalars). All three are rank-1 reductions---a \code{map} over the items the
driver produced, with no cross-coupling between rows.

\begin{keyidea}
The three reduction shapes are: the \textbf{rank-2 symmetric Gram} (capacitance,
inductance---one verb, the weight is the only difference); the \textbf{rank-2
port-projection} (S-parameters---a projection, deliberately not a Gram); and the
\textbf{rank-1 scalar or mode table} (eigenfrequency$+Q$, energy and
participation, waveguide modes---one row per item). A matrix-shaped result does
not by itself mean a Gram: the projection shares the Gram's rank but not its
symmetry or its meaning.
\end{keyidea}

\section{The driver \texorpdfstring{$\times$}{x} framework grid}

We can now lay the six drivers on the framework all at once. \Cref{tab:grid}
is the grid the rest of the book fills in. Each row is a driver; each column is
a stage of the skeleton---the function space it assembles on, what it assembles,
which solve corner its solve sits in, which reduction shape its reduce uses, and
the product it yields. Every cell is a forward reference: the chapters of
Parts~II through~V supply the contents of these cells in full, one stage at a
time.

\begin{table}[t]
\centering
\caption[The driver $\times$ framework grid]{The driver $\times$ framework grid:
the six drivers laid on the skeleton. Reading across a row tells you how that
analysis is built; reading down a column tells you how that stage varies across
analyses. The \emph{solve corner} and \emph{reduction} columns are the two that
this chapter named; almost every difference between rows lives there. The
electrostatic row is the simplest; every other row is the electrostatic row
with one or two stages swapped (last paragraph).}
\label{tab:grid}
\footnotesize
\setlength{\tabcolsep}{3.5pt}
\renewcommand{\arraystretch}{1.3}
\begin{tabular}{@{}llp{27mm}ll@{}}
\toprule
Driver & FE space & \multicolumn{1}{l}{Assemble} & Solve corner & Reduction $\to$ product \\
\midrule
Electrostatic & $\Hone$ & $\Kop$ once
  & fixed-operator map & Gram ($w{=}1$) $\to$ capacitance \\
Magnetostatic & $\Hcurl$ & $\Kop$ once
  & fixed-operator map & Gram ($w{=}\tfrac{1}{I_iI_j}$) $\to$ inductance \\
Driven & $\Hcurl$ & $\Kop,\Cop,\Mop$ once; $\mat{A}(\omega)$ rebuilt in loop
  & operator-varying map & projection $\to$ S-parameters \\
Transient & $\Hcurl$ & $\Kop,\Cop,\Mop$ once
  & state-threaded fold & table $\to$ field history \\
Eigenmode & $\Hcurl$ & $\Kop,\Cop,\Mop$ once
  & opaque black-box & table $\to$ frequencies, $Q$ \\
Boundary mode & $\Hcurl{\oplus}\Hone$ & block pencil (2-D)
  & opaque black-box & table $\to$ waveguide modes \\
\bottomrule
\end{tabular}
\end{table}

Read down the \emph{Solve corner} column and the chapter's central claim becomes
a glance: two drivers map over a fixed operator, one maps over a varying
operator, one folds a threaded state, and two hand the whole problem to a black
box. Read down the \emph{Reduction} column and the three shapes appear: two
Grams, one projection, three tables. The physics columns (space, assemble)
supply variety too, but they are the subject of Part~II; the two columns this
chapter named are the ones that organize the book.

There is a punchline, and it is worth stating plainly. \emph{Electrostatics is
the universal entry point, and every other analysis is electrostatics with one
stage swapped.} Magnetostatics swaps the function space ($\Hone\to\Hcurl$) and
the Gram weight, and changes nothing else. The driven analysis keeps the map but
swaps the solve corner from fixed to operator-varying, and swaps the reduction
from Gram to projection. The transient analysis swaps the map for a fold.
Eigenmode swaps the loop-we-write for a black box. Boundary mode is eigenmode on
a two-dimensional cross-section. Six analyses, one anchor, a handful of swaps:
that is the whole architecture of the simulator, and it is why the book can
teach electrostatics once, in depth (\cref{ch:electrostatics}), and then present
each remaining analysis as a small, well-localized set of differences from it.

\summarysection
The skeleton is
\texttt{config} $\to$ \texttt{assemble} $\to$ \texttt{solve} $\to$
\texttt{reduce} $\to$ \texttt{product},
optionally wrapped in an adaptive-refinement outer fold. Almost all the variety
between the six analyses lives in two stages. The \texttt{solve} stage has four
\emph{corners}: the fixed-operator map (\code{solve\_family}; electrostatic,
magnetostatic), the operator-varying map (\code{frequency\_sweep}; driven), the
state-threaded fold (\code{fold\_solve}; transient and the AMR loop), and the
opaque black-box (\code{eigsolve}; eigenmode, boundary-mode). The \texttt{reduce}
stage has three \emph{shapes}: the rank-2 symmetric Gram (capacitance and
inductance, one verb differing only by weight), the rank-2 port-projection
(S-parameters, a projection and deliberately not a Gram), and the rank-1
scalar-or-mode table (eigenfrequency$+Q$, energy, waveguide modes). The driver
$\times$ framework grid (\cref{tab:grid}) places all six on these axes, and the
punchline organizes the rest of the book: electrostatics is the universal entry
point, and every other analysis swaps one stage.

\furtherreading
The four-corners and three-shapes taxonomy is this book's own organizing
device, distilled from the structure of the Palace solver
itself~\citep{palace2023}; reading the source's six driver entry points side by
side is the best way to see the shared skeleton in the wild. The distinction
between an independent map and a state-threaded fold is the functional
programmer's daily bread; we build it from scratch in \cref{ch:fold}, assuming
nothing. For the numerical content behind the corners---why the fixed-operator
map can amortize a factorization, why the eigenvalue corner is best left to a
specialized library---Saad's account of iterative methods~\citep{saad2003} is
the standard reference, and we draw on it throughout Part~IV.

\exercisessection
\begin{enumerate}[nosep]
  \item For each of the four solve corners, state the two-question
  classification that places it: is the operator fixed or varying, and is the
  computation a map or a fold? (The black-box corner answers ``neither row''---
  explain why.)
  \item The driven analysis (corner~2) and electrostatics (corner~1) both run a
  family of independent solves. Pinpoint the single line of pseudocode whose
  position---inside the loop versus outside it---is the entire difference between
  them. Why does moving it change the \emph{answer}, not merely the speed?
  \item Capacitance and inductance share one reduction verb and differ only in a
  weight. Write down each weight, and explain in one sentence why inductance
  needs the current normalization that capacitance does not.
  \item S-parameters and capacitance both produce a matrix, yet the book insists
  they are different reduction shapes. Give two structural reasons the
  S-parameter matrix is not a Gram. (Hint: consider symmetry, and consider what
  each index ranges over.)
  \item Using \cref{tab:grid}, describe the transient analysis as ``the
  electrostatic analysis with these stages swapped.'' Exactly which cells change,
  and which stay the same?
\end{enumerate}
